%%==========================================================================
%%  saus-thesis-demo.tex
%%  A sample final year project report showing every part of the thesis
%%  class: cover, title page, approval, declaration, abstracts, lists,
%%  chapters, figures, tables, equations, theorems, code and references.
%%  The students and supervisor are placeholders.
%%  Compile with:  latexmk saus-thesis-demo.tex   (LuaLaTeX and Biber)
%%==========================================================================
% degree  = bs (default) | ms | mphil | phd
% type    = project | thesis          (default: project for bs, thesis otherwise)
% body    = serif (default) | sans
% spacing = onehalf (default) | single | double
% bib     = ieee (default) | apa | numeric | authoryear | none
% twoside   print on both sides; chapters open on right-hand pages
% print     black links, for the bound copy
% status  = final (default) | draft   marks every page "Draft" with the date
\documentclass[degree=bs,bib=ieee]{saus-thesis}

\addbibresource{saus-thesis-demo.bib}

%% ---- Details ---------------------------------------------------------------
\title{A Shared Document Identity for The Shaikh Ayaz University}
\subtitle{Templates for slides, posters, letters, certificates and theses}
\student{Student One}{Roll No. 00-CS-001}
\student{Student Two}{Roll No. 00-CS-002}
\student{Student Three}{Roll No. 00-CS-003}
\supervisor{Dr. A. Supervisor}{Assistant Professor, Department of Computer Science}
\discipline{Computer Science}
\sausdepartment{Department of Computer Science}
% \faculty{Faculty of ...}
\session{2022--2026}
% \submissiondate{September 2026}     % defaults to the current month
\similarity{9\%}                      % leave out if not required
\keywords{document design, typography, \LaTeX, accessibility, Sindhi, Urdu}

% Signatures on the certificate of approval, after the supervisor
\approvedby{External Examiner}{Name, designation and institution}
\approvedby{Head of Department}{Department of Computer Science}

\begin{document}

\makecover            % full-page cover for the PDF copy (leave out when printing)
\frontmatter
\maketitlepage
\makeapproval
\makedeclaration
\dedication{To the teachers of Shikarpur.}

\begin{acknowledgements}
We thank our supervisor for guidance throughout this project, the IT Directorate
for access to the university's crest and photographs, and our classmates who
tested the templates in their own coursework.
\end{acknowledgements}

\begin{abstract}
Documents produced across a university---lecture slides, research posters,
official letters, certificates and theses---are often designed separately, and
so present the institution inconsistently. This project builds a single
document identity for The Shaikh Ayaz University, Shikarpur: a shared brand
layer that holds the colours sampled from the university crest, the typefaces
and the institutional details, and a family of \LaTeX{} templates that draw on
it. The templates set the university's name in English, Sindhi and Urdu, meet
the WCAG~2.1 contrast requirements for text, and compile with LuaLaTeX,
XeLaTeX and pdfLaTeX. Matching PowerPoint, Word and Excel templates extend the
same identity to colleagues who do not use \LaTeX{}.
\end{abstract}

\begin{urduabstract}
یہ منصوبہ شیخ ایاز یونیورسٹی، شکارپور کے لیے دستاویزات کے مشترکہ سانچوں کا ایک
مجموعہ پیش کرتا ہے۔ ان سانچوں میں سلائیڈز، پوسٹر، سرکاری خطوط، اسناد اور مقالے
شامل ہیں، جو یونیورسٹی کے نشان، رنگوں اور حروف کو یکساں انداز میں استعمال کرتے ہیں۔
یونیورسٹی کا نام انگریزی، سندھی اور اردو تینوں زبانوں میں درج کیا جاتا ہے۔
\end{urduabstract}

\tableofcontents
\listoffigures
\listoftables
\lstlistoflistings

\begin{abbreviations}
  \abbr{FYP}{Final Year Project}
  \abbr{HEC}{Higher Education Commission}
  \abbr{PDF}{Portable Document Format}
  \abbr{RTL}{Right to left}
  \abbr{SAUS}{The Shaikh Ayaz University, Shikarpur}
  \abbr{WCAG}{Web Content Accessibility Guidelines}
\end{abbreviations}

\mainmatter

%% ---------------------------------------------------------------------------
\chapter{Introduction}\label{ch:intro}

A university speaks through its documents. A student meets it first in an
admission letter, a colleague elsewhere in a conference poster, and an employer
in a degree certificate. When each of these is designed separately, the
institution appears as several different organisations.

\section{Problem statement}

The Shaikh Ayaz University, Shikarpur, a public sector general university with
two faculties and seven degree programmes, had no shared templates for its
documents. Departments prepared slides, posters and letters independently, with
differing colours, typefaces and renderings of the university's name.

\section{Objectives}

The project set out to:
\begin{enumerate}
  \item derive a colour palette and typographic system from the university crest;
  \item build templates for slides, posters, letters, certificates and theses
        that share this system by construction, not by copying;
  \item set the university's name correctly in Sindhi and Urdu as well as English;
  \item keep every text colour legible under WCAG~2.1~\cite{wcag21}.
\end{enumerate}

\section{Scope}

The templates target \LaTeX{}~\cite{lamport1994}, the standard for technical
writing in the sciences, with matching Office templates for colleagues who
work in PowerPoint, Word and Excel. Chapter~\ref{ch:design} describes the
design, Chapter~\ref{ch:implementation} its implementation and
Chapter~\ref{ch:evaluation} the evaluation.

%% ---------------------------------------------------------------------------
\chapter{Background}\label{ch:background}

\section{Typography and identity}

Bringhurst describes typography as the craft of endowing language with a
durable visual form~\cite{bringhurst2012}. For an institution, that form is
part of its identity: the same typefaces, used the same way, make documents
recognisably its own.

\section{\TeX{} and \LaTeX{}}

\TeX{}~\cite{knuth1984} separates the content of a document from its design.
\LaTeX{} classes and packages carry the design, so a single change to a class
restyles every document built on it. KOMA-Script~\cite{kohm-koma} provides the
foundations for the letter and thesis classes, TikZ~\cite{tantau-tikz} draws the
title pages, and \texttt{fontspec}~\cite{robertson-fontspec} loads the typefaces.

\section{Arabic-script names}

Sindhi and Urdu are written right to left in the Arabic script, Sindhi usually
in Naskh and Urdu in Nastaliq. The templates set Sindhi in Lateef~\cite{sil-lateef}
and Urdu in Noto Nastaliq Urdu, both of which are bundled with the templates.

%% ---------------------------------------------------------------------------
\chapter{Design}\label{ch:design}

\section{The crest}

The university crest (Figure~\ref{fig:crest}) supplies the two primary colours:
the maroon of its script and border, and the royal blue of its outer ring.

\begin{figure}[htbp]
  \centering
  \includegraphics[height=45mm]{\sauslogo}
  \caption{The university crest, the source of the primary colours.}
  \label{fig:crest}
\end{figure}

\section{Colour and legibility}

\begin{definition}[Contrast ratio]\label{def:contrast}
For two colours with relative luminances $L_1 \ge L_2$, the contrast ratio is
\begin{equation}\label{eq:contrast}
  C = \frac{L_1 + 0.05}{L_2 + 0.05},
\end{equation}
which runs from $1$ (no contrast) to $21$ (black on white)~\cite{wcag21}.
\end{definition}

WCAG~2.1 asks for $C \ge 4.5$ for body text and $C \ge 3$ for large text.
Table~\ref{tab:palette} lists the palette with each colour's contrast
against white, computed with Equation~\eqref{eq:contrast}.

\begin{table}[htbp]
  \centering
  \caption{The brand palette and its contrast against white.}
  \label{tab:palette}
  \begin{tabular}{@{}llrl@{}}
    \toprule
    \textbf{Colour} & \textbf{Hex} & \textbf{Contrast} & \textbf{Use} \\
    \midrule
    Maroon      & \texttt{\#6B2A4C} & 10.2:1 & Primary \\
    Deep maroon & \texttt{\#4A1C34} & 13.9:1 & Large fields \\
    Royal blue  & \texttt{\#1A1A91} & 13.1:1 & Secondary \\
    Gold        & \texttt{\#C39B3E} &  2.6:1 & Rules and accents only \\
    Ink         & \texttt{\#231C21} & 16.7:1 & Body text \\
    Slate       & \texttt{\#6A626B} &  5.9:1 & Secondary text \\
    \bottomrule
  \end{tabular}
\end{table}

Gold falls below the threshold on white, so the templates use it only for rules
and ornaments, or for small capitals on the deep maroon field.

%% ---------------------------------------------------------------------------
\chapter{Implementation}\label{ch:implementation}

\section{A shared brand layer}

Every template loads one package, \texttt{saus-brand.sty}, which defines the
colours, typefaces and identity (Figure~\ref{fig:architecture}). A change to
that file reaches every template.

\begin{figure}[htbp]
  \centering
  \begin{tikzpicture}[
      box/.style={draw=sausmaroon,fill=sausmist,rounded corners=2pt,
        minimum width=24mm,minimum height=9mm,font=\sffamily\small},
      brand/.style={box,fill=sausmaroon,text=white,minimum width=40mm},
      >=stealth]
    \node[brand] (brand) {saus-brand.sty};
    \foreach \name/\x in {Slides/-54mm,Posters/-27mm,Letters/0mm,Certificates/27mm,Theses/54mm}
      \node[box] (\name) at (\x,-22mm) {\name};
    \foreach \name in {Slides,Posters,Letters,Certificates,Theses}
      \draw[->,sausmaroonlight,thick] (brand.south) -- (\name.north);
  \end{tikzpicture}
  \caption{Every template draws on the shared brand layer.}
  \label{fig:architecture}
\end{figure}

\section{Using a template}

A presentation needs only the theme; the university's details are already
built in. Listing~\ref{lst:slides} shows a complete title slide.

\begin{lstlisting}[language={[LaTeX]TeX},caption={A title slide with the SAUS theme.},label={lst:slides}]
\documentclass[aspectratio=169]{beamer}
\usetheme{SAUS}                  % colours, fonts and crest
\title{Digital Transformation of Campus Services}
\author{Dr. A. Presenter}
\begin{document}
\begin{frame}[plain]\titlepage\end{frame}
\end{document}
\end{lstlisting}

\section{Engines}

\begin{proposition}\label{prop:engines}
Each template compiles without errors under LuaLaTeX, XeLaTeX and pdfLaTeX.
\end{proposition}

LuaLaTeX and XeLaTeX use the OpenType typefaces directly and set Sindhi and
Urdu; pdfLaTeX falls back to Type~1 versions of the typefaces and omits the
Arabic-script names, as Chapter~\ref{ch:evaluation} confirms.

%% ---------------------------------------------------------------------------
\chapter{Evaluation}\label{ch:evaluation}

\section{Compilation}

Each template's demonstration file was compiled with all three engines.
Table~\ref{tab:engines} summarises the result.

\begin{table}[htbp]
  \centering
  \caption{Compilation of the demonstration files.}
  \label{tab:engines}
  \begin{tabular}{@{}lccc@{}}
    \toprule
    \textbf{Template} & \textbf{LuaLaTeX} & \textbf{XeLaTeX} & \textbf{pdfLaTeX} \\
    \midrule
    Slides       & Yes & Yes & Yes\textsuperscript{a} \\
    Posters      & Yes & Yes & Yes\textsuperscript{a} \\
    Letters      & Yes & Yes & Yes\textsuperscript{a} \\
    Certificates & Yes & Yes & Yes\textsuperscript{a} \\
    \bottomrule
    \multicolumn{4}{@{}l}{\footnotesize\textsuperscript{a}\,Without the Sindhi and Urdu names.}
  \end{tabular}
\end{table}

\section{Legibility}

Every colour used for text reaches a contrast of at least $4.5$ against white
(Table~\ref{tab:palette}), satisfying Definition~\ref{def:contrast} at the
level WCAG~2.1 requires for body text.

%% ---------------------------------------------------------------------------
\chapter{Conclusion}\label{ch:conclusion}

The project produced a shared document identity for the university and a
family of templates that apply it consistently. Because the design lives in one
brand layer, future changes---a new typeface, a revised address---need to be
made only once. Future work could extend the templates to the university's
website and to forms.

\printreferences

\appendix
\chapter{Using the templates}\label{app:using}

The templates, with starter files and full documentation, are in the
\texttt{latex} folder. Copy the template and its starter file, with
\texttt{saus-brand.sty} and the \texttt{assets} folder, into a new directory
and build with \texttt{latexmk}.

\end{document}
